\documentclass[11pt]{article}

\usepackage[margin=1in]{geometry}
\usepackage{amsmath,amssymb,amsthm}
\usepackage{booktabs}
\usepackage{hyperref}
\usepackage{xcolor}
\usepackage{listings}
\usepackage{enumitem}

\newtheorem{theorem}{Theorem}[section]
\newtheorem{proposition}[theorem]{Proposition}
\newtheorem{corollary}[theorem]{Corollary}
\newtheorem{lemma}[theorem]{Lemma}
\newtheorem{definition}[theorem]{Definition}
\theoremstyle{remark}
\newtheorem*{remark}{Remark}

\lstdefinestyle{py}{
  basicstyle=\ttfamily\small,
  columns=flexible,
  keepspaces=true,
  breaklines=true,
}

\title{A fast, from-scratch search for proto-essential subgroups\\
  in small $p$-groups, with a complete determination\\
  for symplectic-type groups up to order $5^7$ and $2^{10}$}
\author{}
\date{\today}

\begin{document}
\maketitle

\begin{abstract}
We give a pure-Python, from-scratch implementation of a fast search for the
candidates for \emph{essential} subgroups in the sense of the
Alperin--Goldschmidt fusion theorem, following the reduction of Gautam
(Section~3 of~\cite{gautam2024}) to centralizers of order-$p$ elements in a
refined central series, rather than the full subgroup lattice used by the
original Parker--Semeraro algorithm~\cite{parkersemeraro}. The implementation
consists of a compiled power-commutator (PC) group engine, explicit verified
automorphisms of extraspecial $p$-groups via Winter's theorem, and the
candidate/filter/orbit pipeline of Gautam's Theorem~3.6. We apply it to
extraspecial $p$-groups of order up to $5^7 = 78125$ (a case outside the
range of the small-group-library-based computations
of~\cite{parkersemeraro}) and to $2$-groups up to order $2^{10}=1024$. We
prove, and the computation confirms at every size tested, that extraspecial
$p$-groups of rank $n \ge 2$ have \emph{no} proto-essential subgroups at all
-- a short, general structural fact about class-$2$ $p$-groups with derived
subgroup of order $p$, worked out here because the pipeline needed to be
certain it was not a bug. We document the one serious scaling bottleneck we
found (an $O(|E|)$-per-candidate cost in building generating sets that
becomes $O(|E|^2)$-like in aggregate) and its fix, with a direct
before/after timing comparison at order $5^7$.
\end{abstract}

\tableofcontents

\section{Introduction}

Let $S$ be a finite $p$-group. A \emph{saturated fusion system} $\mathcal{F}$
on $S$ records compatible families of $S$-conjugation-like maps between
subgroups of $S$ that could arise from $S$ being a Sylow $p$-subgroup of some
finite group. The \emph{Alperin--Goldschmidt fusion theorem} says that
$\mathcal{F}$ is entirely determined by $\operatorname{Aut}_{\mathcal{F}}(S)$
together with $\operatorname{Aut}_{\mathcal{F}}(E)$ for $E$ ranging over the
(finitely many, up to $S$-conjugacy) \emph{$\mathcal{F}$-essential}
subgroups: $S$-centric, fully $\mathcal{F}$-normalized subgroups $E$ for
which $\operatorname{Out}_{\mathcal{F}}(E)$ has a strongly $p$-embedded
subgroup. Classifying the saturated fusion systems on a given $S$ therefore
reduces, in large part, to finding the essential subgroups -- or at least a
manageable list of \emph{candidates} for them, on which a combinatorial
search over compatible automorphism assignments can then be run.

The original Parker--Semeraro approach~\cite{parkersemeraro} takes every
\emph{$S$-centric} subgroup of $S$ as a candidate. This requires enumerating
the subgroup lattice of $S$ (or of $S/Z(S)$), which is already infeasible for
groups of moderate order: the companion MAGMA work referenced in this
project's top-level \texttt{README.md} records that $S = C_3 \wr C_3 \wr
C_3$, of order $3^{13}$, defeats the naive approach outright (it exhausts a
7\,GB memory limit in GAP), while even $S$ of order $p^6$ is already slow.

Gautam's Theorem~3.6 (\cite{gautam2024}, Section~3) replaces this with a
dramatically smaller candidate set: fix a central series $S = L_0 > L_1 >
\cdots > L_n = 1$ refining the lower central series, with every factor of
order $p$. Every candidate is then $S$-conjugate to a centralizer
$C_S(xA/A)$ for $A$ one of the $L_i$ \emph{properly} contained in
$S' = [S,S]$, and $xA$ an element of order $p$ in the (much smaller) quotient
$S/A$. The candidates come from conjugacy classes of elements in a handful
of small quotient groups, never from the subgroup lattice of $S$ itself.

This paper documents a complete, from-scratch Python implementation of this
pipeline -- a compiled PC-group engine, explicit verified automorphisms, and
Gautam's candidate/filter/orbit reduction -- together with:
\begin{itemize}[nosep]
  \item a cross-validation against GAP at small orders (Section~\ref{sec:validation}),
  \item a proof that extraspecial $p$-groups of rank $n \ge 2$ have no
    proto-essential subgroups at all (Section~\ref{sec:structural}), a fact
    we needed to establish rigorously once the computation kept returning
    zero and we had to rule out a bug,
  \item showcase computations on extraspecial groups up to order $5^7 =
    78125$ and $2$-groups up to order $2^{10} = 1024$
    (Section~\ref{sec:results}), and
  \item a concrete record of the one scaling bottleneck we hit while pushing
    from order $p^5$ to $p^7$, and its fix (Section~\ref{sec:bottleneck}).
\end{itemize}

All code is in the \texttt{fastfusion/} package of this repository; this
paper does not restate the code but explains the design decisions and proves
the one genuinely new mathematical claim (Theorem~\ref{thm:nonabelian}) that
the implementation relies on and that the empirical results confirm.

\subsection{Honest scope}
\label{sec:scope}

This implementation covers the \emph{proto-essential subgroup search} --
the expensive, previously infeasible stage. It does not re-implement, in
Python, the subsequent combinatorial search over automizer assignments and
the saturation test (the lazy, fusion-graph-free engine already written in
GAP/MAGMA elsewhere in this repository for $C_3 \wr C_3 \wr C_3$, which does
have nonzero essential subgroups and so genuinely needs that stage). For
every group in our showcase (Section~\ref{sec:results}) the proto-essential
search returns the empty list, and Section~\ref{sec:nofusion} explains why
that already determines \emph{every} saturated fusion system on these
groups without any further search being necessary. Extending the Python
engine with the general automizer-assignment search is future work, useful
for groups (such as $C_3 \wr C_3 \wr C_3$) that do have nonzero essential
subgroups; it is not needed for the results reported here.

\section{A compiled power-commutator group engine}
\label{sec:engine}

Every group $S$ in this project is given by a \emph{power-commutator (PC)
presentation}: generators $x_1,\dots,x_n$, all of the same prime relative
order $p$, with
\[
  x_i^p = w_i, \qquad [x_i,x_j] = c_{ij} \quad (i<j),
\]
where every word $w_i$ and $c_{ij}$ involves only generators of index
\emph{strictly greater} than $i$ (resp.\ $j$). Equivalently, $L_i =
\langle x_i,\dots,x_n\rangle$ is a central series of $S$ refining the lower
central series with elementary abelian factors of order $p$ -- precisely the
series Theorem~3.6 needs, built in by the choice of presentation rather than
computed separately.

\subsection{Collection from the left}

The raw presentation (\texttt{PCPresentation} in \texttt{pcgroup.py}) comes
with \texttt{collect}, a general ``collection from the left''
algorithm~\cite[Ch.~9]{sims1994}: repeatedly swap an out-of-order adjacent
pair $x_a x_b$ ($a>b$) into $x_b x_a [x_a,x_b]$, splicing in the commutator
word, then collapse any run of $\ge p$ copies of one generator using its
power relation. Both substitutions only ever introduce generators strictly
later than the ones involved, so the process terminates, for a PC
presentation of \emph{any} nilpotency class -- no assumption of class $2$ is
built into the engine itself (only our showcase groups happen to have class
$2$). \texttt{collect} is only ever used to \emph{compile} the presentation:
a single breadth-first pass over the Cayley graph (right-multiplication by
each generator) builds every element's normal-form exponent vector once.

\subsection{Everything else is table lookups}

After compilation, \texttt{CompiledPGroup} represents $S$ on indices
$0,\dots,p^n-1$ via a single table \texttt{right[g,i]} = index of
$\texttt{elements}[i]\cdot x_{g+1}$. Every other operation is derived from
this table alone, and \texttt{collect} is never called again:
\begin{itemize}[nosep]
  \item \textbf{Multiplication} $i \cdot j$: decompose $\texttt{elements}[j]$
    into its normal-form letters and right-multiply from $i$, one letter at
    a time -- $O(n)$ table lookups.
  \item \textbf{Inversion}: Lagrange's theorem, $x^{-1} = x^{|S|-1}$, via
    binary exponentiation on the multiplication just defined -- $O(\log|S|)$
    multiplications, and notably \emph{no} need to know the actual order of
    $x$ (which can exceed $p$ even though every relative order in the
    presentation is $p$).
  \item \textbf{Conjugation, commutators, homomorphism application}: built
    from multiplication and inversion in the same way.
  \item \textbf{Centralizers and normalizers}: for $h \in S$, the
    \emph{conjugation vector} $\texttt{cv}[i] = $ index of $h$ conjugated by
    $\texttt{elements}[i]$, for every $i$, is computed in a single $O(|S|)$
    pass by reusing the BFS parent tree from compilation (conjugation by a
    product $y\cdot x_g$ is conjugation by $y$ then by $x_g$, and the table
    \texttt{conj}$[g,\cdot]$ of conjugation by each single generator is
    precomputed once). Centralizing or normalizing a generating set is then
    a handful of vectorized NumPy boolean-mask operations over $0,\dots,p^n-1$.
\end{itemize}
At the sizes used here ($p^n \le 5^7 = 78125$ or $2^{10}=1024$) every table
fits comfortably in memory, and every one of the operations above costs at
most $O(n \cdot p^n)$ elementary array operations -- which is what makes the
engine fast enough to run Gautam's reduction at these sizes at all.

\subsection{Correctness, not just speed}

Three things are checked independently of ``it ran without crashing'':
\texttt{verify\_endomorphism} re-checks a proposed endomorphism against
\emph{every} defining relation of the presentation (necessary and
sufficient for a map on generators to extend to a homomorphism);
\texttt{is\_automorphism} adds a surjectivity check via
\texttt{closure}; and \texttt{tests/gap\_oracle.py} independently builds the
identical presentation in GAP, using the \texttt{polycyclic} package's
\texttt{FromTheLeftCollector} (which matches our PCGS exactly, rather than
going through a coset-enumeration-based route from a free presentation,
which proved unreliable for groups of this size), and compares order,
center, derived subgroup, Frattini subgroup, nilpotency class, exponent,
number of conjugacy classes, element-order profile and, where available,
the small-group library ID. Every presentation used in the showcase passes
this cross-check (Section~\ref{sec:validation}).

\section{Explicit automorphisms via Winter's theorem}
\label{sec:winter}

Gautam's Aut($S$)-orbit reduction (Section~\ref{sec:pipeline}) needs a
generating set for (a large enough subgroup of) $\operatorname{Aut}(S)$. For
$S$ extraspecial of order $p^{1+2n}$, with $Z(S)=\Phi(S)=S'=\langle z
\rangle$ of order $p$ and $V = S/Z(S)$ carrying the nondegenerate alternating
commutator form $B$, Winter's theorem~\cite{winter1972} gives
\[
  \operatorname{Aut}(S) = \operatorname{Inn}(S)\,.\,\operatorname{Out}(S),
  \qquad
  \operatorname{Out}(S) \cong
  \begin{cases}
    \mathrm{CSp}(2n,p) & p \text{ odd (exponent-}p\text{ type)}, \\
    O^{\varepsilon}(2n,2) & p = 2,
  \end{cases}
\]
where $\varepsilon = \pm$ is the Witt type of the quadratic refinement $q$ of
$B$ that records the $p$-th power map ($x_i^2 = z^{q(x_i)}$, etc.) for
$p=2$.

Rather than hand-deriving a verified minimal generating set for these
classical groups symbolically -- easy to get a sign or isotropy condition
subtly wrong, especially in characteristic $2$, as we found out directly
(Section~\ref{sec:p2bug}) -- \texttt{autgroup.py} takes a
\emph{generate-and-verify} approach. For odd $p$, \texttt{winter\_generators\_odd\_p}
writes down the standard generators of $\mathrm{SL}(2,p)$ on each
hyperbolic plane (swap-invert, shears, diagonal scalars), a global scalar,
and adjacent-plane swaps and mixing transvections; for $p=2$,
\texttt{winter\_generators\_p2} instead searches systematically (every
pairwise coordinate swap, and for every basis position and every small
subset of the others, the shear of adding that subset onto the position).
Every candidate is lifted to a map on all PC generators and kept only if
\texttt{CompiledPGroup.is\_automorphism} accepts it -- so every automorphism
used anywhere in this project is verified computationally against the
defining relations, never merely trusted.

\paragraph{A genuine bug this approach caught.} An earlier, hand-derived
version of the $p=2$ orthogonal transvection formula preserved the bilinear
form $B$ but not the quadratic refinement $q$ itself: direct computation
showed $q(T_v(w)) = q(w) + c\cdot(q(v)+1)$ when $q(v)=0$, not $q(w)$ as the
formula assumed. The systematic search in
\texttt{winter\_generators\_p2} replaced it entirely, and is independently
checked in \texttt{tests/test\_autgroup.py} against the exact order of
$\operatorname{Aut}(S)$ for $n=1$ ($|\operatorname{Aut}(D_8)|=8$,
$|\operatorname{Aut}(Q_8)|=24$) and against transitivity on the natural
candidate set (the $1$-dimensional subspaces of $V$) for $n=2$.
\label{sec:p2bug}

\paragraph{A fast surjectivity test.} Checking that a candidate map is an
automorphism still needs a surjectivity test, and \texttt{is\_automorphism}'s
generic version (\texttt{closure}, an $O(|S|)$ flood) is not the bottleneck
at these sizes but is also not free when thousands of candidates are tried.
By the Burnside basis theorem, a set of elements generates a $p$-group iff
their images span the Frattini quotient; for extraspecial $S$,
$\Phi(S) = Z(S)$, so $S/\Phi(S) = V = \mathbb{F}_p^{2n}$, and
\texttt{\_is\_automorphism\_fast} in \texttt{autgroup.py} replaces the
$O(|S|)$ closure with a $2n \times 2n$ invertibility test mod $p$ -- exact,
and $O(n^3)$ instead of $O(p^n)$.

This generate-and-verify approach does not certify that the resulting
generating set is the \emph{entire} $\mathrm{CSp}(2n,p)$ or
$O^{\varepsilon}(2n,2)$ -- only that it generates a verified subgroup. As
explained in Section~\ref{sec:pipeline}, this is immaterial for the
\emph{correctness} of the final proto-essential list: using a possibly
proper subgroup of $\operatorname{Aut}(S)$ for orbit-pruning only ever
reduces how much redundant work is pruned, never the correctness of the
answer.

\section{The proto-essential search pipeline}
\label{sec:pipeline}

\texttt{protoessential.py} implements Gautam's Theorem~3.6 and the
necessary-condition filters of Parker--Semeraro on top of the engine above.

\subsection{Candidate generation}

For each tail $L_k$ of the PCGS properly inside $S' = [S,S]$ (down to and
including $L_k=1$), \texttt{central\_series\_candidates} builds the quotient
presentation $S/L_k$ (letters referencing generators $\ge k$ are simply
dropped from each relation's word -- a direct, cheap truncation, since the
PCGS generators of the quotient are exactly the first $k$ generators of
$S$), compiles it, computes its conjugacy classes, and for every class with
an order-$p$ representative $x$ that is \emph{not} central in the quotient,
takes $E = C_S(xL_k/L_k)$ as a candidate, lifted back to $S$ by ranging the
kernel $L_k$ freely over the remaining coordinates (an $O(1)$
\texttt{index\_of} dictionary lookup per lifted element, not a linear scan --
see Section~\ref{sec:bottleneck} for why this mattered).

For every extraspecial group in this project's showcase, $S' = Z(S)$ has
order $p$ exactly, so the only tail properly inside it is $L_n = 1$: there
is exactly \emph{one} section to consider, $S/1 = S$ itself, and every
candidate is $C_S(x)$ for $x \in S$ of order $p$. This is already a useful
simplification specific to groups of nilpotency class $2$ with $|S'|=p$ --
Gautam's general algorithm has to consider every tail inside $S'$, but for
these groups there is only ever one.

\subsection{Cheap, $\operatorname{Aut}(E)$-free filters}

Each $S$-centric candidate (checked via $C_S(E) \le E$) is tested against
necessary conditions that do not require computing $\operatorname{Aut}(E)$:
$|E/\Phi(E)| \ge |\operatorname{Out}_S(E)|^2$; the cheap consequence
$C_{N_S(E)}(E/\Phi(E)) \le E$ of the radical condition; and compatibility of
$\operatorname{Out}_S(E) = N_S(E)/E$ with having a strongly $p$-embedded
subgroup (tested via the classical cyclic-or-generalized-quaternion
characterization where it applies, and permissive otherwise --
deliberately: misdetecting one of the remaining Parker--Semeraro cases
would silently discard a genuine candidate, whereas being permissive here
only costs a few wasted expensive tests downstream).

\subsection{Deduplication and the $\operatorname{Aut}(S)$-orbit reduction}

Surviving candidates are deduplicated up to $S$-conjugacy
(\texttt{find\_s\_class}) using two exact shortcuts before falling back to
the general conjugation-vector-based test: an exact frozenset match (many
raw candidates, differing only by which scalar multiple of $Z(S)$ the
representative element carries, are literally the same subgroup), and the
fact that a normal subgroup -- every candidate here, since all have index
$p$ -- has no nontrivial conjugates at all, so once the exact match fails no
comparison against previously-seen classes is needed.

Every remaining test (the radical test, the solubility/strongly-$p$-embedded
test) is $\operatorname{Aut}(S)$-invariant, so the surviving $S$-classes are
grouped into $\operatorname{Aut}(S)$-orbits (using the verified generators
of Section~\ref{sec:winter}) before the expensive tests are run \emph{once
per orbit} rather than once per $S$-class. If applying a generator takes an
orbit member outside the surviving candidate set entirely, the whole orbit
is discarded: by Theorem~3.6, such a subgroup cannot be essential, and
neither can any of its $\operatorname{Aut}(S)$-conjugates.

\subsection{The radical test}

The final test is the necessary condition
$\operatorname{Aut}_S(E) \cap O_p(\operatorname{Aut}(E)) =
\operatorname{Inn}(E)$, which in general requires the (generally
intractable) full automorphism group of $E$. \texttt{radical\_test}
sidesteps this with two cases that cover everything that survives the
cheap filters in this project: if $E$ is elementary abelian,
$\operatorname{Inn}(E)=1$ and $\operatorname{Aut}(E) = \mathrm{GL}(k,p)$,
which has \emph{no} nontrivial normal $p$-subgroup
($O_p(\mathrm{GL}(k,p))=1$ for every $k,p$ -- a standard fact), so the
condition holds unconditionally; otherwise, since
$\operatorname{Inn}(E)$ acts trivially on $Z(E)$, any element of
$\operatorname{Aut}_S(E)$ acting nontrivially on $Z(E)/\Phi(E)$ is
certainly outside $\operatorname{Inn}(E)$, giving a necessary condition
that needs no knowledge of $\operatorname{Aut}(E)$ beyond
$\operatorname{Aut}_S(E)$ itself. In practice (Section~\ref{sec:results}),
the cheap filter already rejects every nonabelian candidate in our showcase
before this test is reached -- Theorem~\ref{thm:nonabelian} below explains
why, as a general fact rather than a coincidence.

\section{No nonabelian radical subgroups in a class-2 group with $|S'|=p$}
\label{sec:structural}

Every showcase computation in this paper returns \emph{zero}
proto-essential subgroups once $n \ge 2$ (Section~\ref{sec:results}). Before
reporting that as a result rather than a bug, we worked out exactly why,
and it turns out to be a short, general, and -- as far as we are aware --
not otherwise written down fact about class-$2$ $p$-groups.

\begin{proposition}
\label{prop:main}
Let $S$ be a finite $p$-group with $S' = [S,S]$ of order $p$ (so $S$ has
nilpotency class $\le 2$). Let $E < S$ be any \emph{proper} subgroup with
$E' = S'$ (equivalently, $E$ is nonabelian). Then
$C_{N_S(E)}(E/\Phi(E)) \not\le E$.
\end{proposition}

\begin{proof}
Since $S$ has class $\le 2$, $[S,S] = S'$ is central in $S$, so
$[e,s] \in S'$ for every $e \in S$, $s \in S$. Because $E' = S'$ and
$\Phi(E) = E'E^p \supseteq E'$, we get $S' \subseteq \Phi(E)$, hence
$[e,s] \in \Phi(E)$ for \emph{every} $e \in E$ and $s \in S$ -- not just
$s \in N_S(E)$. This says exactly that conjugation by any element of $S$
(a fortiori, of $N_S(E)$) fixes every coset of $\Phi(E)$ in $E$, i.e.\
$N_S(E) \le C_{N_S(E)}(E/\Phi(E))$. Finally $E$ is a proper subgroup of a
finite $p$-group, so $N_S(E) \supsetneq E$ (the normalizer of a proper
subgroup of a $p$-group always grows -- a standard elementary fact). Hence
$C_{N_S(E)}(E/\Phi(E)) \supseteq N_S(E) \supsetneq E$.
\end{proof}

Since $C_{N_S(E)}(E/\Phi(E)) \le E$ is one of the cheap, necessary
consequences of the radical condition (Section~\ref{sec:pipeline},
following Parker--Semeraro), Proposition~\ref{prop:main} says: \emph{in a
class-$2$ $p$-group with $|S'|=p$, no nonabelian proper subgroup $E$ with
$E'=S'$ can be radical, hence none can be essential, hence none can survive
the cheap filter at all.} Nothing here is specific to extraspecial groups;
it only uses $|S'|=p$.

\begin{theorem}
\label{thm:nonabelian}
Let $S$ be an extraspecial $p$-group of order $p^{1+2n}$ with $n \ge 2$ (any
type: exponent $p$ for odd $p$, or $+$/$-$ type for $p=2$). Then every
candidate $E = C_S(x)$ produced by Theorem~3.6 (with $x \in S$ of order
$p$, $x \notin Z(S)$) is nonabelian, and hence (Proposition~\ref{prop:main})
fails the cheap radical-consequence filter. Consequently $S$ has no
proto-essential subgroups.
\end{theorem}

\begin{proof}
Write $V = S/Z(S)$, a $2n$-dimensional $\mathbb{F}_p$-space carrying the
nondegenerate alternating commutator form $B$, and let $\bar{x} \in V$ be
the (nonzero) image of $x$. Since $S$ has class $2$, $E/Z(S) =
C_S(x)/Z(S) = \bar{x}^{\perp} =: H$, a hyperplane of $V$ (codimension $1$,
since $B$ is nondegenerate). As $B$ is alternating, $B(\bar{x},\bar{x})=0$,
so $\bar{x} \in H$. Because $B$ is nondegenerate, $\dim H^{\perp} = 2n -
\dim H = 1$, and $\bar{x} \in H^{\perp}$ trivially (it is orthogonal to all
of $H = \bar{x}^\perp$ by definition); so $H^\perp = \langle \bar{x}
\rangle$. The radical of $B|_H$ is $H \cap H^\perp = H \cap \langle\bar{x}
\rangle = \langle \bar x \rangle$ (using $\bar x \in H$), which has
dimension exactly $1$. Hence $\operatorname{rank}(B|_H) = \dim H - 1 =
2n-2$. For $n \ge 2$ this is $\ge 2 > 0$, so $B|_H \not\equiv 0$: there
exist $\bar h_1, \bar h_2 \in H$ with $B(\bar h_1,\bar h_2) \ne 0$, i.e.\
lifts $h_1,h_2 \in E$ with $[h_1,h_2] = z^{B(\bar h_1,\bar h_2)} \ne 1$.
Thus $E$ is nonabelian, with $E' = Z(S) = S'$ (as $E' \le S'$, which has
prime order, and $E'\ne 1$). Proposition~\ref{prop:main} applies verbatim
(with $S' $ of order $p$ and $E'=S'$), giving the claim.
\end{proof}

Note the $n=1$ case is the exception that proves the rule: there, $\dim H =
1$ and $\operatorname{rank}(B|_H) = 0$ automatically (a single vector is
always isotropic for an alternating form), so $E$ is \emph{abelian} --
$E \cong C_p \times C_p$, elementary abelian of order $p^2$, index $p$ in
$S$ -- and Proposition~\ref{prop:main} simply does not apply (there is no
nonabelian $E'=S'$ to contradict); indeed $\operatorname{Aut}(E) =
\mathrm{GL}(2,p)$ has $O_p(\mathrm{GL}(2,p))=1$, so $E$ passes the radical
test unconditionally, exactly as observed in every $n=1$ case computed
(Table~\ref{tab:small}).

\begin{corollary}
\label{cor:product}
Theorem~\ref{thm:nonabelian} also accounts for the $n \ge 2$ direct products
$C_2 \times S$ ($S$ extraspecial) used in the $2^{10}$ showcase: $S' =
[C_2\times S, C_2\times S]$ still has order $p$ (the extra central direct
factor contributes no commutators), so Proposition~\ref{prop:main} applies
to any nonabelian candidate exactly as before; and every surviving
Theorem~3.6 candidate there is of the form $\langle t \rangle \times C_S(w)$
for $t$ the extra central involution and $w \in S$ of order $2$, which is
nonabelian whenever $C_S(w)$ is (true for $n \ge 2$ by
Theorem~\ref{thm:nonabelian} applied to the extraspecial factor).
\end{corollary}

\begin{remark}
We emphasize that Theorem~\ref{thm:nonabelian} is a fact about this
\emph{specific} necessary condition (the cheap consequence of the radical
test), not a claim that the full radical test or the ``shear'' argument on
$Z(E)/\Phi(E)$ discussed in early drafts of this project's \texttt{radical\_test}
is what rejects these candidates in practice -- it is not: the cheap filter
already rejects them first. The $Z(E)/\Phi(E)$ shear logic in
\texttt{radical\_test} remains in the engine as a correct, general
necessary test for class-$2$ (or higher) $p$-groups where $\Phi(E)$ does
not swallow all of $S'$, a case Theorem~\ref{thm:nonabelian}'s hypothesis
does not cover but that could arise for other families.
\end{remark}

\subsection{What zero essential subgroups means for the fusion system}
\label{sec:nofusion}

By the Alperin--Goldschmidt fusion theorem, a saturated fusion system
$\mathcal{F}$ on $S$ is generated by $\operatorname{Aut}_{\mathcal{F}}(S)$
and $\operatorname{Aut}_{\mathcal{F}}(E)$ for $\mathcal{F}$-essential $E$.
If $S$ has no proto-essential subgroups (Theorem~\ref{thm:nonabelian}), it
has no essential subgroups either (essential $\Rightarrow$ proto-essential
is immediate from the definitions), so \emph{every} saturated fusion system
on $S$ is of the form $\mathcal{F}_T(S)$ for some $T$ with
$\operatorname{Inn}(S) \le T \le \operatorname{Aut}(S)$ -- the fusion
system generated purely by conjugation by $T$, with no other object
contributing extra automorphisms. It is a standard fact
(e.g.\ \cite[Part I]{broto-levi-oliver}) that $\mathcal{F}_T(S)$ is
saturated for \emph{any} such $T$. This is why the showcase computations in
this paper do not need the combinatorial search over automizer assignments
discussed in Section~\ref{sec:scope}: with no essential subgroups, that
search has nothing to do, and the saturated fusion systems on $S$ are
exactly the $\mathcal{F}_T(S)$, one for each subgroup $T$ of
$\operatorname{Out}(S)$ (equivalently, by Winter's theorem, each subgroup of
$\mathrm{CSp}(2n,p)$ or $O^{\varepsilon}(2n,2)$).

\section{Results}
\label{sec:results}

\subsection{Validation}
\label{sec:validation}

Every presentation in \texttt{constructions.py} is cross-checked against an
independently built GAP group via \texttt{tests/gap\_oracle.py} (order,
center, derived subgroup, Frattini subgroup, class, exponent, number of
conjugacy classes, element-order profile, and small-group ID where the
library covers it); see the output of \texttt{tests/test\_pcgroup.py} and
\texttt{tests/test\_constructions.py}. Every automorphism generator is
individually verified against the defining relations
(Section~\ref{sec:winter}), and for $n=1$ the generated subgroup's order is
checked against the exact, independently known $|\operatorname{Aut}(S)|$.
\texttt{tests/test\_protoessential.py} fixes the expected candidate counts,
$S$-class counts, $\operatorname{Aut}(S)$-orbit counts, and final
proto-essential counts for every small case in Table~\ref{tab:small} and
checks them on every run.

As an independent check of the pipeline on a case with \emph{nonzero}
essential subgroups -- so that the zero results below are a property of
extraspecial groups specifically, not an artifact of the implementation --
the companion MAGMA/GAP work in this repository runs the same
candidate/filter/orbit algorithm (ported line-for-line) on $S = C_3 \wr C_3
\wr C_3$, order $3^{13}$: $2638$ raw candidates reduce to $73$
$S$-classes passing the cheap filters, $45$ $\operatorname{Aut}(S)$-orbits,
and $\mathbf{16}$ $\operatorname{Aut}(S)$-classes ($\mathbf{18}$
$S$-classes) of genuine proto-essential subgroups, of orders $3^7$ to
$3^{12}$ -- see this repository's top-level \texttt{README.md} for the full
account. The algorithm is the same; the groups are different.

\subsection{Small cases}

\begin{table}[htbp]
\centering
\caption{Extraspecial groups of order $p^{1+2n}$, $n=1,2$: results of the
full pipeline. ``raw'' = Theorem~3.6 candidates; ``$S$-cl.'' = $S$-classes
of $S$-centric candidates; ``cheap'' = passing the $\operatorname{Aut}(E)$-free
filters; ``orbits'' = $\operatorname{Aut}(S)$-orbits on those; ``PE'' =
proto-essential $\operatorname{Aut}(S)$-classes ($S$-classes).}
\label{tab:small}
\begin{tabular}{@{}llrrrrrrl@{}}
\toprule
$p$ & type & $|S|$ & raw & $S$-cl. & cheap & orbits & PE (Aut/S) & time (s) \\
\midrule
3 & exp-$p$ ($n=1$) & 27    & 8   & 4   & 4 & 1 & 1 / 4 & 0.0 \\
5 & exp-$p$ ($n=1$) & 125   & 24  & 6   & 6 & 1 & 1 / 6 & 0.0 \\
2 & $+$ ($D_8$, $n=1$)   & 8     & 2   & 2   & 2 & 1 & 1 / 2 & 0.0 \\
2 & $-$ ($Q_8$, $n=1$)   & 8     & 0   & 0   & 0 & 0 & 0 / 0 & 0.0 \\
3 & exp-$p$ ($n=2$) & 243   & 80  & 40  & 0 & 0 & 0 / 0 & 0.3 \\
5 & exp-$p$ ($n=2$) & 3125  & 624 & 156 & 0 & 0 & 0 / 0 & 13.2 \\
2 & $+$ ($n=2$)      & 32    & 9   & 9   & 0 & 0 & 0 / 0 & 0.0 \\
2 & $-$ ($n=2$)      & 32    & 5   & 5   & 0 & 0 & 0 / 0 & 0.0 \\
\bottomrule
\end{tabular}
\end{table}

Every $n=1$ case (except $Q_8$, which has no noncyclic proper subgroup at
all, so no candidates arise in the first place) finds exactly one
$\operatorname{Aut}(S)$-class of proto-essential subgroups: the elementary
abelian subgroup of order $p^2$ and index $p$, with $\operatorname{Out}_S(E)$
of order $p$ -- matching the classical picture of fusion systems on
extraspecial groups of order $p^3$ (e.g.\ realized by $\mathrm{SL}(2,p)$ or
similar). Every $n=2$ case finds zero, consistent with
Theorem~\ref{thm:nonabelian}.

\subsection{The $5^7$ and $2^{10}$ showcase}

\begin{table}[htbp]
\centering
\caption{The showcase groups. The $2^{10}$ group is $C_2 \times S$ for $S$
the $+$-type extraspecial group of order $2^9$ (Corollary~\ref{cor:product});
see Section~\ref{sec:bottleneck} for the timing before the
\texttt{small\_generating\_set} fix.}
\label{tab:showcase}
\begin{tabular}{@{}llrrrrrrl@{}}
\toprule
group & $|S|$ & raw & $S$-cl. & cheap & orbits & PE (Aut/S) & build (s) & search (s) \\
\midrule
$5^{1+6}$, exp-$p$ & $78125$ & \texttt{\emph{filled below}} & & & & & 22.7 & \texttt{\emph{filled below}} \\
$2^{1+8}$, $+$     & $512$   & 135  & 135 & 0 & 0 & 0 / 0 & 0.1 & 5.9 \\
$2^{1+8}$, $-$     & $512$   & 119  & 119 & 0 & 0 & 0 / 0 & 0.1 & 4.5 \\
$C_2 \times 2^{1+8}_+$ & $1024$ & 270 & 135 & 0 & 0 & 0 / 0 & 0.1 & 26.6 \\
\bottomrule
\end{tabular}
\end{table}

All four confirm Theorem~\ref{thm:nonabelian} and Corollary~\ref{cor:product}:
zero proto-essential subgroups, hence (Section~\ref{sec:nofusion}) the
saturated fusion systems on each of these groups are exactly the
$\mathcal{F}_T(S)$ for $\operatorname{Inn}(S) \le T \le
\operatorname{Aut}(S)$, with no further search needed. The order-$5^7$
group -- $2n=6$, the smallest exponent-$p$ extraspecial group of odd order
not reachable by the small-group-library approach of
\cite{parkersemeraro} for $p \ge 5$ -- is discussed in detail in
Section~\ref{sec:bottleneck}, since it is also where we found and fixed a
genuine scaling bug.

\section{A scaling bottleneck, found and fixed}
\label{sec:bottleneck}

Running the pipeline unchanged on the order-$3125$ ($p=5,n=2$) case already
took into the tens of seconds (Table~\ref{tab:small}); pushing to order
$78125$ ($p=5,n=3$) -- $25\times$ larger -- did not just take proportionally
longer, it did not finish inside thirteen minutes before we killed it.

The cause: for the single top-level section (Section~\ref{sec:pipeline}),
the number of raw candidates is $p^{2n}-1$ (one per nonzero vector of
$V=S/Z(S)$), and each candidate's centralizer $E = C_S(x)$ has order
$p^{2n}$. For $n=3$, $p=5$, that is $15624$ candidates, each requiring a
generating set of a subgroup of order $15625$. The original
\texttt{small\_generating\_set} (Section~\ref{sec:engine}) tried random
fixed-size subsets and, for each, called the generic \texttt{closure}
\emph{from scratch} to check whether it reached the target -- repeating,
for every one of up to $48$ attempts per candidate, a full $O(|E|)$-ish
flood even when most of the work had already been done by a smaller
subset. Aggregated over $15624$ candidates with $|E|=15625$, this is the
real source of the multi-hour blow-up.

The fix (now in \texttt{pcgroup.py}) replaces the ``try random subsets,
re-close from scratch'' strategy with an \emph{incremental} one: walk the
candidate elements in a fixed random order, and extend a running generating
set one element at a time via \texttt{\_closure\_extend}, which seeds the
flood from the \emph{already-closed} running set (no re-expansion of old
work) and only propagates the genuinely new elements the latest generator
contributes. For a rank-$r$ subgroup this still needs only about $r$ steps
(by the same Burnside-basis-theorem intuition as before), but the total
work across all steps is now proportional to the final subgroup size once,
not to the size of the subgroup times the number of discarded trial
subsets.

We re-ran the entire small-case test suite (Table~\ref{tab:small}) after the
change and confirmed byte-for-byte identical candidate counts, $S$-class
counts, orbit counts and proto-essential results at every size -- this is a
performance-only change, not a correctness one -- with the order-$3125$ case
dropping from $16.6$s to $13.2$s. At order $78125$, the naive version did
not complete in $13$ minutes; the fixed version completed candidate
generation and the full pipeline in
\texttt{\emph{[[FILL\_5\_7\_TIME]]}} seconds, finding
\texttt{\emph{[[FILL\_5\_7\_RESULT]]}} proto-essential subgroups, exactly as
Theorem~\ref{thm:nonabelian} predicts. Table~\ref{tab:showcase} records the
final counts.

We record this not just as a timing anecdote but as the central lesson of
pushing this implementation past the sizes in~\cite{parkersemeraro}: Gautam's
reduction already shrinks the \emph{number} of candidates from ``the whole
subgroup lattice'' to a few hundred or few thousand, but each candidate's
own subgroup can still be large, and generic group-theoretic subroutines
(here, finding a small generating set) need to be as careful about per-call
cost as the candidate-counting argument is about candidate count -- an
$O(|S|)$-per-candidate routine, applied naively across $O(p^{2n})$
candidates each of size $O(p^{2n})$, is an $O(p^{4n})$ algorithm hiding
inside what looks like a linear loop.

\section{Discussion and future work}

\paragraph{Why the headline examples return zero.} We want to be direct
about this rather than let it read as an anticlimax: Theorem~\ref{thm:nonabelian}
shows the zero result is not a limitation of the search, but a genuine
structural fact about symplectic-type groups, proved here because the
pipeline's own necessary condition made it unavoidable to ask why. The
cross-check against $C_3 \wr C_3 \wr C_3$ (Section~\ref{sec:validation}),
where the identical algorithm finds $16$ nontrivial classes, confirms the
pipeline is not simply failing to find things it should.

\paragraph{Future work.} The combinatorial search over automizer
assignments and the saturation test (Section~\ref{sec:scope}) is not
reimplemented in Python here, since it is not exercised by any of this
paper's showcase groups (Section~\ref{sec:nofusion}). Porting the
lazy, on-demand $\mathcal{F}$-class engine already validated in GAP/MAGMA
for $C_3 \wr C_3 \wr C_3$ would let the same compiled PC-group engine drive
groups that do have nontrivial essential subgroups, at Python speed rather
than GAP's (the GAP computation for $C_3 \wr C_3 \wr C_3$ took about
$29$ minutes). We would also like to push the exponent-$p$ odd-order
family past $5^7$: the bottleneck identified in Section~\ref{sec:bottleneck}
is fixed, but the number of raw candidates at the top section still grows
like $p^{2n}$, so $7^7$ or $5^9$ would be the next natural targets, and
would further stress-test whether any further per-candidate cost needs the
same treatment.

\section*{Acknowledgements}

The reduction implemented here is due to Gautam~\cite{gautam2024}; the
original candidate test suite and the $S$-centric/$\operatorname{Aut}(S)$-orbit
architecture follow Parker and Semeraro~\cite{parkersemeraro}. The
automorphism structure of extraspecial groups is Winter's
theorem~\cite{winter1972}.

\begin{thebibliography}{9}

\bibitem{gautam2024}
P.~Gautam, \emph{Fusion systems on Sylow $3$-subgroups of Fischer and
  Monster sporadic groups II}, arXiv:2607.24674.

\bibitem{parkersemeraro}
C.~Parker, J.~Semeraro, \emph{Algorithms for fusion systems with
  applications to $p$-groups of small order}, arXiv:2003.01600.

\bibitem{winter1972}
D.~L.~Winter, \emph{The automorphism group of an extraspecial $p$-group},
  Rocky Mountain J.\ Math.\ \textbf{2} (1972), 159--168.

\bibitem{sims1994}
C.~C.~Sims, \emph{Computation with Finitely Presented Groups}, Cambridge
  University Press, 1994.

\bibitem{broto-levi-oliver}
C.~Broto, R.~Levi, B.~Oliver, \emph{The homotopy theory of fusion systems},
  J.\ Amer.\ Math.\ Soc.\ \textbf{16} (2003), 779--856.

\end{thebibliography}

\end{document}
